% ================================================================
% ==== FILE: content/m_spaces/m5.tex
% ================================================================

% ==============================
%  Ontology of Continua — Core
%  Meta-Space: M5
%  FULL MODULE — FINAL VERSION
% ==============================

\subsubsection{\texorpdfstring{$M_5$}{M_5} Übersicht}
\label{sec:m5-overview}

$M_5$ ist der Metaraum, der \emph{elektrische Erregbarkeit} zulässig macht.
Es ist die Umgebung, in der ein biologisches Kontinuum Folgendes zeigen kann:

\begin{itemize}
    \item stabile Membranspannung $V(t)$,
    \item Ionenkanal-vermittelte Leitfähigkeitszustände,
    \item schwellenbasierte Anregung ($\Theta_{\mathrm{exc}}$),
    \item Ausbreitung elektrischer Ereignisse (Vorläufer von Aktionspotentialen),
    \item regeneratives Feedback für Spitzenzyklen erforderlich.
\end{itemize}

Wo $M_4$ die \emph{chemische} Kompartimentierung ermöglichte,
$M_5$ ermöglicht \emph{elektrochemische} Dynamiken, die das $K_5$-Kontinuum definieren.

% ---------------------------------------------------------------
\subsubsection*{1. Zulässiger Konfigurationsraum \texorpdfstring{$\Omega(M_5)$}{\Omega(M_5)}}

Die zulässigen Konfigurationen erweitern $\Omega(M_4)$ um:

\[
\Omega(M_5) =
\left\{
(\partial\Omega,\ V,\ g_{\mathrm{ion}},\
\Pi_{\mathrm{ion}},\
\Delta\mu_{\mathrm{ion}},\
C_m,\
\mathcal{C}_{\mathrm{Kanäle}})
\ \Big|\
\text{Zulässigkeitsbedingungen gelten}
\right\}.
\]

Wo:

\begin{itemize}
    \item $V$ — Membranspannungsfeld,
    \item $C_m$ — Membrankapazität,
    \item $g_{\mathrm{ion}}$ — Leitfähigkeitszustände von Ionenkanälen,
    \item $\mathcal{C}_{\mathrm{channels}}$ – Kanalpopulation und Gating-Logik,
    \item $\Pi_{\mathrm{ion}}$ — ionenspezifische Permeabilitäten,
    \item $\Delta\mu_{\mathrm{ion}}$ – elektrochemische Gradienten, die Ströme erzeugen.
\end{itemize}

Voraussetzung für die Zulassung sind:

\begin{enumerate}
    \item Membranintegrität wie in $M_4$.
    \item Ionenkanäle, die zwischen mindestens zwei Zuständen wechseln können:
          \[
          \{ \text{offen},\ \text{geschlossen} \}.
          \]
    \item Spannungsdynamik erfüllt:
          \[
          C_m \frac{dV}{dt} = - \sum_i g_i (V - E_i) + I_{\mathrm{ext}}
          \]
          in einem zulässigen Parameterbereich.
    \item Existenz eines nichttrivialen Gleichgewichts und einer Separatrix, die a definiert
          Schwelle.
\end{enumerate}

Somit ist $M_5$ die minimale Umgebung, in der eine Zelle erregbar sein kann.

% ---------------------------------------------------------------
\subsubsection*{2. Zulässige Achsen \texorpdfstring{$A(M_5)$}{A(M_5)}}

Es entsteht eine neue Grundachse:

\[
A_{\mathrm{exc}} = \{ \text{erregbar},\ \text{nicht erregbar} \},
\]

mit der entsprechenden Spannungs-Schwellen-Unterscheidung:

\[
A_V = \{ V < \Theta_{\mathrm{exc}},\ V \ge \Theta_{\mathrm{exc}} \}.
\]

Weitere zulässige Achsen:

\begin{itemize}
    \item $A_{\mathrm{ion}}$ — ionenspezifische Kanalkonformationen,
    \item $A_{\mathrm{gate}}$ – Gating-Variablenzustände (Aktivierung/Inaktivierung),
    \item $A_{\mathrm{cond}}$ – Leitfähigkeitsregime,
    \item $A_{\mathrm{spike}}$ – Refraktär-/Überschwing-/Rücksetzphasen.
\end{itemize}

Diese Achsen definieren die hochdimensionale dynamische Geometrie anregbarer Membranen.

% ---------------------------------------------------------------
\subsubsection*{3. Potentiale \texorpdfstring{$P(M_5)$}{P(M_5)}}

$M_5$ führt mit der Elektrodynamik verbundene Potentiale ein:

\[
P(M_5) =
(U_{\mathrm{ion}},\ U_V,\ U_{\mathrm{gate}},\
U_{\mathrm{CAP}},\ U_{\mathrm{rest}}).
\]

Interpretation:

\begin{itemize}
    \item $U_{\mathrm{ion}}$ — Nernst-ähnliche elektrochemische Potentiale,
    \item $U_V$ — spannungsabhängige Energielandschaft,
    \item $U_{\mathrm{gate}}$ – Gating-Energie für Kanalübergänge,
    \item $U_{\mathrm{CAP}}$ — kapazitive Energie der Membran,
    \item $U_{\mathrm{rest}}$ — Energieminimum im Ruhezustand.
\end{itemize}

Voraussetzungen für die Zulassung:

\begin{enumerate}
    \item $U_V$ muss ein lokales Minimum haben (Ruhezustand),
    \item Die Verstärkungsfunktion muss eine überschwellige Instabilität zulassen,
    \item Gating-Übergänge müssen thermisch und chemisch machbar sein.
\end{enumerate}

% ---------------------------------------------------------------
\subsubsection*{4. Schwellenwerte \texorpdfstring{$\Theta(M_5)$}{\Theta(M_5)}}

Die definierende Schwelle:

\[
\Theta_{\mathrm{exc}} = \text{minimale Depolarisation erforderlich, um ein regeneratives Ereignis auszulösen}.
\]

Zusätzliche Schwellenwerte:

\begin{itemize}
    \item $\Theta_{\mathrm{gate}}$ – Gating-Aktivierungsschwelle,
    \item $\Theta_{\mathrm{cond}}$ — Schaltschwelle der Leitfähigkeit,
    \item $\Theta_{\mathrm{Stabilität}}$ — Stabilitätsgrenze für Schwingungsmodi.
\end{itemize}

Kritische Beziehung:

\[
V(t) \ge \Theta_{\mathrm{exc}} \quad\Rightarrow\quad
\text{Übergang in ein Spike-erzeugendes Regime}.
\]

% ---------------------------------------------------------------
\subsubsection*{5. Flüsse \texorpdfstring{$J(M_5)$}{J(M_5)}}

Zulässige Flüsse erweitern die in $M_4$ um Folgendes:

\[
J(M_5) =
\big(
J_{\mathrm{ion}},\
J_V,\
J_{\mathrm{Gate}},\
J_{\mathrm{exc}},\
\nabla_i T_5
\Big).
\]

Wo:

\begin{itemize}
    \item $J_{\mathrm{ion}}$ — ionenspezifische Ströme,
    \item $J_V$ – spannungsgetriebene Flüsse,
    \item $J_{\mathrm{gate}}$ – dynamische Gating-Übergänge,
    \item $J_{\mathrm{exc}}$ – regeneratives Anregungsflussregime,
    \item $\nabla_i T_5$ – Strömungen, die durch Spannungsgradienten unter Erregbarkeit angetrieben werden.
\end{itemize}

Naturschutzgesetze:

\[
\sum_i J_{\mathrm{ion},i} = C_m \frac{dV}{dt}.
\]

% ---------------------------------------------------------------
\subsubsection*{6. Zyklen in \texorpdfstring{$M_5$}{M_5}}

Der Schlüsselzyklus:

\[
C_{\mathrm{Spike}}
=
\{\text{Rest} \to \text{Depolarisation} \to
\text{Überschwingen} \to \text{Repolarisation} \to \text{Refraktärität} \to \text{Rest}\}.
\]

Weitere zulässige Zyklen:

\begin{itemize}
    \item $C_{\mathrm{gate}}$ – Gating-Aktivierungs-/Inaktivierungsschleife,
    \item $C_{\mathrm{osc}}$ – Oszillationszyklen unterhalb des Schwellenwerts (sofern zulässig),
    \item $C_{\mathrm{conduct}}$ – Übergangszyklen zwischen Leitfähigkeitszuständen.
\end{itemize}

Voraussetzung für die Existenz eines Spike-Zyklus:

\[
\oint_{C_{\mathrm{Spike}}} dA_{\mathrm{Phase}} \ approx 0
\quad\text{und}\quad
V(t) \text{ kreuzt } \Theta_{\mathrm{exc}} \text{ einmal pro Zyklus}.
\]

% ---------------------------------------------------------------
\subsubsection*{7. Grenze \texorpdfstring{$\partial\Omega(M_5)$}{\partial\Omega(M_5)}}

Eine Konfiguration nähert sich $\partial\Omega(M_5)$, wenn:

\begin{itemize}
    \item die Membran kann kein stabiles Ruhepotential unterstützen,
    \item Gating-Kinetik unterbricht die Erregbarkeit,
    \item Leitfähigkeit sättigt oder bricht zusammen,
    \item $\Theta_{\mathrm{exc}}$ wird unerreichbar (zu hoch) oder trivial (null),
    \item Depolarisationsblock bleibt bestehen,
    \item Ionengradienten fallen unter minimale Lebensfähigkeitsniveaus.
\end{itemize}

Das Überschreiten von $\partial\Omega(M_5)$ macht eine Spitze unmöglich;
somit bricht $K_5$ zusammen oder kehrt zu einem $K_4$-ähnlichen Regime zurück.

% ---------------------------------------------------------------
\subsubsection*{8. Beziehung zu \texorpdfstring{$M_4$}{M_4} und \texorpdfstring{$M_6$}{M_6}}

Der Übergang $K_4 \to K_5$ ist zulässig, wenn:

\[
A_{\mathrm{exc}} \in A(M_5),
\qquad
V(t)\ge\Theta_{\mathrm{exc}},
\qquad
\text{Ionenkanäle bilden kohärente Gating-Dynamik}.
\]

Der Übergang zu $K_6$ (kognitive Attraktoren) erfordert:

\[
A_{\mathrm{Konzept}} \in A(M_6),
\qquad
C_{\mathrm{Spike}} \text{ muss die Kopplung mit anderen Spikes unterstützen},
\]
\[
\text{Entstehung der Attraktordynamik aus Spike-Zügen}.
\]

Somit ist $M_5$ der Metaraum, der Folgendes ermöglicht:

\begin{itemize}
    \item elektrische Erregbarkeit,
    \item Spitzengenerierung,
    \item dynamische Leitfähigkeitsregulierung,
    \item frühe bioelektrische Logik,
    \item die strukturelle Brücke zur neuronalen Berechnung.
\end{itemize}
